%قبل از شروع مسئله ترجیحاً چیزی ننویسید. اما اگر خواستید بنویسید، در قالب لیستی (فقرات و شمارش) بنویسید.
%برای اضافه کردن نمره، بعدی از شروع‌مسئله، [x] را اضافه کنید و به جای x  نمره مورد نظر را قرار دهید. برای زیرمسئله‌ها هم به همین صورت است.
%\شروع‌مسئله[10]
\شروع‌مسئله

% صورت اصلی مسئله را در این‌جا بنویسید.
فراز که در گراف اوستا شده است، آوازه‌ش به عنوان یک گراف‌شناس معروف در جهان به گوشه همه رسیده است. به همین دلیل، از او دعوت شده است که در شهری عجیب، یک سیمنار علمی برگزار کند. در این شهر، 
$n$
میدان وجود دارد و بین برخی از میدان‌ها، جاده‌هایی وجود دارد. می‌توان به میدان‌ها به شکل رئوس یک گراف نگاه کرد و جاد‌ه‌ها، یال‌هایی بین رئوس هستند. در این شهر، بین هر دو میدانی، دقیقا یک مسیر وجود دارد. ضمنا جاده‌ها، طول‌های متفاوتی می‌توانند داشته باشند. حال فراز که می‌خواهد نشان بدهد به فرهنگ این شهر اهمیت می‌دهد، می‌خواهد بلندترین مسیر ساده بین دو میدان در این شهر را بیابد. راهی از اردر 
$O(n)$
به او بدهید که بتواند به این مهم، دست یابد.
\solution{
	برای قسمت اول، فرض کنید که $Q_k$ احتمال قرار گرفتن دقیقا $k$ کلید در یک جایگاه مشخص باشد. تعداد راه‌های انتخاب این $k$ کلید به صورت 
	$\binom{n}{k}$
	است. آن‌ها با احتمال
	$\frac{1}{n}$
	به همان جایگاه مشخص نگاشت می‌شوند و $n-k$ کلید باقی‌مانده هم هر یک با احتمال
	$1 - \frac{1}{n}$
	به یکی از $n-1$ جایگاه دیگر نگاشت می‌شوند. در نتیجه:
	
	$$Q_{k}=\left(\frac{1}{n}\right)^{k}\left(\frac{n-1}{n}\right)^{n-k}\left(\begin{array}{l}n \\ k\end{array}\right)$$
	
	برای قسمت دوم، می‌دانیم که  حد بالای احتمال این که جایگاه با بیش‌ترین تعداد کلید دقیقا $k$ کلید داشته باشد، برابر وجود جایگاه‌ای با دقیقا $k$ کلید است. $n$ جایگاه وجود دارد که می‌توانند این $k$ کلید را در خود جای دهند. در نتیجه
	$P_k \leq n Q_k$
	
	برای قسمت سوم، با توجه به تقریب استرلینگ می‌توانیم نتیجه بگیریم که
	$\binom{n}{k} \leq (\frac{ n e}{k})^k$
	در نتیجه:
	$Q_{k} \leq \frac{(n-1)^{n-k}}{n^{n}}\left(\frac{n e}{k}\right)^{k} \leq e^{k} / k^{k}$
	
	برای قسمت چهارم در اصل هدف م این است که به
	 $Q_{k_0}< n^{-3}$
	 برسیم. از فرمول هر دو طرف $\log$ می گیریم و سعی‌ می‌کنیم در لگاریتمشان این فرمول را ایجاد کنیم. با لگاریتم گرفتن داریم:
	 
	 $$
	 \frac{c \lg n}{\lg \lg n}(\lg e+\lg \lg \lg n-\lg c-\lg \lg n)<-3 \lg n
	 $$
	 
	 با تقسیم دو طرف بر 
	 $- \log \log n$
	 داریم
	 $$
	 3<-\frac{c \lg e}{\lg \lg n}-\frac{\lg \lg \lg n}{\lg \lg n}+\frac{\lg c}{\lg \lg n}+c
	 $$
	 
	 در $n \to \infty$ همه جملات صفر می‌شوند به جز جمله $c$. در نتیجه با انتخاب $c=4$ خواسته سوال به ازای $n>N$ محقق خواهد شد. با توجه به نتیجه قسمت دوم، می‌توانیم براساس نزولی بودن $Q_k$‌ برحسب $k$ نتیجه بگیریم که به ازای $k\geq k_0$ مقدار $P_k < 1/n^2$ است.
	 
	 برای قسمت آخر، می‌توانیم از طریق حالت بندی به عبارت مد نظر برسیم. در اصل می‌توانیم این طور در نظر بگیریم که برای $M$ های بزرگتر از $c \lg n / \lg \lg n$ مقدار $M$ را در حد بالا $n$ در نظر گرفته و برای مقادیر کوچکتر از آن برابر خود $c \lg n / \lg \lg n$ در نظر بگیریم. بدین ترتیب یک حد بالا برای امیدریاضی $M$ خواهیم داشت. چون امید ریاضی به صورت حاصلضرب احتمال در مقدار متغیر تصادفی است و ما در هر دو حالت حد بالایی را برای مقدار متغیر تصادفی در نظر گرفته ایم. با توجه به شرط قسمت قبلی می دانیم که
	 $n P(M > \frac{c \lg n}{\lg \lg n}) <1$
	 و همچنین می‌توان با توجه به ذات احتمالاتی نتیجه گرفت که
	 $\sum_{i=1}^{d} P(M=i) \leq 1$
	 به ازای $d = c \lg n / \lg \lg n$  در نتیجه با جایگذاری به راحتی به این نتیجه می‌رسیم که
	 $E[M] \in O(\lg n / \lg \lg n)$
	

}



%%%%%%%%%%%%%%%%%%%%%%%%%%%%%%%%%%%%%%%%%%%%%%%%%%%%%
% برخی دستورات مفید برای تمرین‌ها. %
%%%%%%%%%%%%%%%%%%%%%%%%%%%%%%%%%%%%%%%%%%%%%%%%%%%%%

% فرمول‌های نمونه:
%$f(x) = \Theta(x)$
%$T(x, y) \leq \log(x) + 2^y + \binom{2n}{n}$
% {\tt code\_function}


% لیست موارد بدون شماره‌گذاری:
%\شروع{فقرات}
%    \فقره مورد اول 
%        \شروع{فقرات}
%            \فقره مورد اول داخلی 
%                \شروع{فقرات}
%                    \فقره اولین مورد خیلی داخلی 
%                    \فقره دومین مورد خیلی داخلی 
%                \پایان{فقرات}
%            \فقره مورد دوم داخلی 
%        \پایان{فقرات}
%    \فقره مورد دوم 
%\پایان{فقرات}

% لیست موارد شماره‌گذاری شده:
%\شروع{شمارش}
%    \فقره یک مورد 
%    \فقره یک مورد دیگر 
%    \فقره باز هم یک مورد دیگر
%\پایان{شمارش}

% برای نوشتن یک فرمول چند خطی که یک شماره بخوره. برای حالت‌های دیگر به سایت زیر مراجعه کنید:
% https://www.overleaf.com/learn/latex/Aligning_equations_with_amsmath
%\begin{align} 
%    \begin{split}
%        (x+y)^3 &= (x+y)^2(x+y) \\
%                &= (x^2+2xy+y^2)(x+y) \\
%                &= (x^3+2x^2y+xy^2) + (x^2y+2xy^2+y^3) \\
%                &= x^3+3x^2y+3xy^2+y^3
%    \end{split}                                 
%\end{align}

% نوشتن ماتریس و آرایه دوبعدی
%\begin{align}
%    A = 
%    \begin{bmatrix}
%        A_{11} & A_{21} \\
%        A_{21} & A_{22}
%    \end{bmatrix}
%\end{align}

% اضافه کردن تصویر با شرح. برای اضافه کردن تصویر بدون شرح، فقط دستور درج‌تصویر را استفاده کنید.
%\شروع{شکل}[h!]
%\تنظیم‌ازوسط
%  \درج‌تصویر[پهنا=0.5\پهنای‌متن]{img.jpg}
%\شرح{شرح تصویر}
%\برچسب{تصویر:۱}
%\پایان{شکل}
%\پایان{زیرمسائل}

%کد زیر را ببینید.
%%\lstinputlisting{source_filename.py}
%\begin{latin}
%\begin{lstlisting}
%# Import the modules
%import sys
%import random
%
%ans = True
%
%while ans:
%question = raw_input("Ask the magic 8 ball a question: (press enter to quit) ")
%
%answers = random.randint(1,8)
%
%if question == "":
%sys.exit()
%
%elif answers == 1:
%print "It is certain"
%
%elif answers == 2:
%print "Outlook good"
%
%elif answers == 3:
%print "You may rely on it"
%
%elif answers == 4:
%print "Ask again later"
%
%elif answers == 5:
%print "Concentrate and ask again"
%
%elif answers == 6:
%print "Reply hazy, try again"
%
%elif answers == 7:
%print "My reply is no"
%
%elif answers == 8:
%print "My sources say no"
%\end{lstlisting}
%\end{latin}
